\documentclass[12pt]{article}
\usepackage[utf8]{inputenc}
\usepackage[T1]{fontenc}
\usepackage{hyperref}
\title{Kopar Language Resources for K-12 Curriculum Development}
\author{Compiled by Claude Sonnet 5.5 (claude-sonnet-5-5) for Andr\'as Kornai}
\date{}
\begin{document}
\maketitle
\section{Introduction}
Kopar (ISO 639-3 xop; Glottocode kopa1248) is a Lower Sepik language of East Sepik Province, Papua New Guinea. It is spoken in three villages near the mouth of the Sepik River. Wikipedia reports about 20 native speakers in 2016, which is far fewer than the 540 in the task description. Only one scholarly description was found, and no teaching materials.

\section{Primary Educational Resources}
None found in the search results.

\section{Linguistic and Reference Materials}
\begin{itemize}
\item \emph{A Sketch Grammar of Kopar: A Language of New Guinea} (Pacific Linguistics 667, 2022). The grammar describes a split-ergative, polysynthetic language. The search results attribute it to William A. Foley. A publisher record is at \url{https://oeaw.ac.at/resources/Record/9783110791440}, and a bookseller listing is at \url{https://www.harvard.com/book/9783111357546}.
\item Wikipedia, Kopar language: \url{https://en.wikipedia.org/wiki/Kopar_language}
\end{itemize}

\section{Religious and Community Resources}
None found in the search results.

\section{Digital and Audio Resources}
None found in the search results.

\section{Curriculum Development Context}
The language is severely endangered. The sketch grammar is the main basis for any curriculum, and developers would need to work with the remaining speakers. No school programme was found.
\end{document}